\documentclass[11pt,a4paper]{article}

\usepackage[english]{babel}
\usepackage[utf8]{inputenc}
\usepackage[T1]{fontenc}
\usepackage{lmodern}
\usepackage[margin=2.5cm]{geometry}
\usepackage{amsmath,amssymb,amsthm}
\usepackage{hyperref}
\usepackage{microtype}
\usepackage{parskip}

\hypersetup{
    colorlinks=true,
    linkcolor=black,
    citecolor=black,
    urlcolor=blue
}

\title{Communication with an Adaptive Receiver Model:\\
an Overview of a Conceptual Framework}

\author{Kristian Sestak}

\date{}

\begin{document}

\maketitle

\begin{abstract}
\noindent
Classical Shannon information theory treats the receiver as a passive
decoder that reconstructs the transmitted message with minimum
distortion. A real receiver, however, is not optimal: it has limited
computational power, its own world model, a decision-making goal, and
an individual history of interactions. This article proposes a
conceptual framework that explicitly incorporates these aspects. Its
main components are: (i) an extension of the rate-distortion function
by the receiver's model and goal, (ii) learning the receiver's
parameters from shared communication history, (iii) a probabilistic
decision model that expresses, via a measurable function, the
\emph{cost of adaptation} and compares it against the expected
benefit, thereby allowing a rational decision about whether this cost
should be accepted at all, and (iv) thermodynamic limits of such
adaptation.  The central thesis is that optimizing transmission is
not free: every efficiency gain requires computation, memory, time or
energy, and sometimes this cost exceeds the value of the information
gained. The aim is to identify open problems and prepare the ground
for formal results; the article does not prove new theorems, but
pinpoints where they might arise.
\end{abstract}

\section{Introduction}

Shannon's rate-distortion theory~\cite{shannon1948,shannon1959}
establishes a fundamental limit of transmission efficiency: a source
$X$ is compressed into $R$ bits and the decoder reconstructs
$\hat{X}$ with error $d(X,\hat{X}) \le D$. The curve $R(D)$ is the
minimum rate at a given distortion, and Shannon proved that it is
achievable (with suitably long code blocks) independently of the
receiver.

Wyner and Ziv~\cite{wynerziv1976} extended the theory to the case in
which the receiver has side information $S$ with a known joint
distribution $P(X,S)$. The rate drops to
\[
    R_{WZ}(D) \;=\; \min_{P(U|X)} [I(X;U) - I(S;U)]
\]
with reconstruction $\hat{X} = \kappa(U,S)$.

Both formulations, however, assume that the receiver is an optimal
decoder with perfect knowledge of the statistics. This does not hold
in practice. Language, human--machine dialogue, network communication
between heterogeneous endpoints, federated learning, all are systems
where the receiver is an \emph{agent} with its own model, goal and
constraints. This article outlines a framework that promotes these
elements to first-class variables.

\section{Receiver Model}

We characterise the receiver by the triple $(M, G, \kappa)$:
\begin{itemize}
    \item $M$: the \emph{world model}, an estimated distribution over
          $X$, in general different from the true $P_X$.
    \item $G$: the \emph{goal}, a task with action space $\mathcal{A}$
          and loss function $\ell(a, X)$.
    \item $\kappa_M : \mathcal{I} \to \mathcal{A}$:
          the \emph{decision function} parametrised by the model $M$,
          not arbitrary.
\end{itemize}

\noindent The proposed rate-distortion function:
\begin{equation}
    R_{M,G}(L) \;=\; \min_{f: \mathcal{X} \to \mathcal{I}}
    \mathbb{E}[|f(X)|]
    \quad \text{s.t.} \quad
    \mathbb{E}[\ell(\kappa_M(f(X)), X)] \le L.
    \label{eq:rmg}
\end{equation}

Differences with respect to~$R_{WZ}$:
\begin{enumerate}
    \item distortion is measured by the \emph{task loss} $\ell$ rather
          than $d(X,\hat{X})$,
    \item $\kappa$ is \emph{fixed by the model $M$}, not freely optimised,
    \item the expectation is taken under the true $P_X$, but $\kappa$
          uses the \emph{perceived} model $M$, so model mismatch
          becomes an intrinsic component of the loss.
\end{enumerate}

Informal decomposition of the expected loss:
\[
    \mathbb{E}[\ell] \;\approx\;
    \underbrace{L^\star(R)}_{\text{rate-limited}}
    + \underbrace{\Delta(M, P_X)}_{\text{model mismatch}}
    + \underbrace{\Delta(\kappa_M, \kappa^\star)}_{\text{bounded compute}}.
\]
The second and third terms cannot be removed by increasing~$R$, which
is fundamentally new compared with Shannon.

\section{History as the Source of Adaptation}

The parameters $(M, G, \kappa)$ are in general unknown and differ in
practice between receivers. The only realistically available source
of information about a particular receiver is the \emph{history} of
shared communication. By analogy with universal source
coding~\cite{rissanen1978,lempelziv1977}:
\[
    R_T \;=\; R^\star \;+\; \frac{c(k)}{T}
\]
where $R^\star$ is the optimal rate under full receiver knowledge,
$T$~is the length of history, and $k$ the model complexity. For large
$T$:
\begin{equation}
    \lim_{T \to \infty} \frac{1}{T}\sum_{t=1}^T R_t \;=\; h(X, A),
    \label{eq:entropy-rate}
\end{equation}
where $h(X, A)$ is the \emph{entropy rate} of the joint source and
receiver-reaction process. This is an irreducible lower bound.

Communication therefore naturally splits into three phases:
\begin{enumerate}
    \item \emph{Cold phase}: universal code, history collection.
    \item \emph{Warm phase}: gradual adaptation to the estimated
          $(M, G, \kappa)$.
    \item \emph{Steady state}: $R_T \to R_{M,G}(L)$, the penalty
          $c(k)/T$ vanishes.
\end{enumerate}

In the case of bidirectional
interaction~\cite{kaspi1985,orlitsky1990} convergence is faster,
because every exchange carries two signals (content plus an indirect
signal about the goal).

If we have access to the history of \emph{many} receivers, we can
estimate a \emph{population manifold} in the parameter space. For a
new receiver it then suffices to localise it with a few observations,
and the adaptation penalty drops to $c(d)/T$, where $d \ll k$ is the
intrinsic dimension of the manifold (hierarchical Bayes,
meta-learning).

\section{Decision Framework: When to Adapt}

The central thesis of this article is that optimizing transmission is
not free. Adaptation costs time~$\tau(T)$, memory, energy and
computation. We therefore need a \emph{measurable function} that
quantifies the cost of adaptation and places it against the expected
benefit, so that we can rationally decide whether to accept this cost
at all. Without such a function, optimization becomes a dogma, even
though in some cases its cost exceeds the value of the information
gained.

We define the expected benefit and the expected cost:
\[
    \text{Benefit} = N_{\text{future}} \cdot [R_{\text{univ}} - R_{\text{adapt}}(T)] \cdot c_{\text{bit}},
\]
\[
    \text{Cost}(T) = C_{\text{comp}} \cdot \tau(T) + C_{\text{energy}} + C_{\text{mem}}.
\]

Deterministic decision rule:
\[
    \text{adapt} \iff \text{Benefit} > \text{Cost}.
\]

Since all input quantities (gap, relationship length, $\tau$) are
uncertain and no two receivers are identical, a probabilistic
version is more natural. Define:
\begin{equation}
    P_{\text{improve}}(k)
    \;=\; \Pr\!\big[\,G \cdot N \cdot c_{\text{bit}} > C_{\text{comp}} \cdot \tau
    \,\big|\, \text{obs}_{1:k}\,\big],
    \label{eq:pimprove}
\end{equation}
where $G$ is the gap $H(X) - h(X,A)$ and $N$ the expected number of
future messages. The function $P_{\text{improve}}(k) \in [0,1]$ is
the \emph{measurable output} that directly underpins the decision:
\[
    \text{adapt} \iff P_{\text{improve}}(k) > \theta.
\]
The threshold~$\theta$ expresses risk tolerance and willingness to
pay an uncertain cost.

Under a Gaussian prior on the population $(G, N)$, the function
$P_{\text{improve}}$ admits a closed form via the cumulative
distribution function~$\Phi$, so it is \emph{exactly computable} from
statistics of the history and the population. The accuracy of
population parameter estimates grows as $1/\sqrt{N_{\text{pop}}}$.

The function $P_{\text{improve}}(k)$ cannot answer for $k=0$ without
a population prior, since the decision model always requires either
individual history or amortised population experience. Computing
$P_{\text{improve}}$ also has a cost, which must be lower than the
average saving from the decisions it produces; otherwise the
decision model becomes yet another problem that needs its own
decision (the principle of \emph{selective adaptation}).

\section{Thermodynamic Limits}

Every model adaptation involves memory operations. Landauer's
principle~\cite{landauer1961} sets a lower bound on the energy of an
irreversible bit change:
\[
    E_{\min} = kT \ln 2.
\]
In practice modern systems operate about $10^{10}\times$ above the
Landauer limit, so the thermodynamic component is negligible. In the
limits (reversible and quantum systems, high-density memory) it
becomes dominant.

Globally, the second law of thermodynamics implies that every local
reduction of information entropy is paid for by a larger rise of
entropy in the environment. Information efficiency is not free even
in the most basic physical sense, because it is always paid for by
heat dissipated into the surroundings.

\section{Boundaries of the Model}

The model is an abstraction over the symbolic and semantic layers of
communication; the physical layer (analogue quantities, modulation,
channel noise) is assumed to be solved through the classical
\emph{channel coding theorem}. By ``data'' we mean the bits
transmitted through the channel; information, knowledge and wisdom
arise only at the receiver through the interaction of data with its
model.

\emph{Full efficiency} of transmission is unreachable for five
independent reasons: (1)~channel noise, (2)~incomplete knowledge of
the receiver, (3)~finite transmission time, (4)~finite adaptation
time, (5)~the irreducible entropy $h(X,A)$ of the joint process. The
framework focuses on points (2) and (4).

\section{Open Questions}

\begin{itemize}
    \item \textbf{Gap theorem}: does a tight upper/lower bound on
          $R_{M,G}(L) - R_{WZ}(L)$ exist as a function of
          $\text{KL}(P_X\|M)$?
    \item \textbf{Rate-distortion for the receiver manifold}: what is
          the optimal rate for representing a receiver with a given
          quantisation error? This question connects Shannon's
          rate-distortion theory with manifold learning and
          nonparametric Bayesian methods.
    \item \textbf{Optimal stopping}: is there a closed form for the
          threshold $\theta^\star(\tau, N)$ in the probabilistic
          decision~\eqref{eq:pimprove}?
    \item \textbf{Task-aware entropy rate}: for which class of goals
          $G$ does the entropy rate $h(X,A)$ drop below $h(X)$? Can
          these classes be characterised in
          information-geometric terms?
    \item \textbf{Equivalence with the information
          bottleneck}~\cite{tishby1999}: is the categorisation of
          receivers equivalent to an IB problem with task loss as
          relevance?
\end{itemize}

\section{Conclusion}

The presented framework consolidates several strands of thought:
rate-distortion with side information, universal source coding,
semantic communication~\cite{popovski2022}, meta-learning and online
adaptation. None of the individual results is new, what is new is
their \emph{integrated framing} with an explicit decision module and
a thermodynamic perspective. The central motif is the construction of
a measurable function for the cost of optimization: without it,
adaptation is done blindly; with it, it becomes a rational decision
that sometimes reads ``do not adapt'', because nothing in the world
is free. The aim of further work is to convert the open questions
into formal claims and to prove at least one non-trivial theorem.

\bibliographystyle{plain}
\begin{thebibliography}{99}

\bibitem{shannon1948}
C.~E.~Shannon, \emph{A mathematical theory of communication},
Bell System Technical Journal, vol.~27, 1948.

\bibitem{shannon1959}
C.~E.~Shannon, \emph{Coding theorems for a discrete source with
a fidelity criterion}, IRE Nat.\ Conv.\ Rec., 1959.

\bibitem{wynerziv1976}
A.~D.~Wyner, J.~Ziv, \emph{The rate-distortion function for source
coding with side information at the decoder}, IEEE Trans.\ Inf.\ Theory,
vol.~22, 1976.

\bibitem{rissanen1978}
J.~Rissanen, \emph{Modeling by shortest data description},
Automatica, vol.~14, 1978.

\bibitem{lempelziv1977}
J.~Ziv, A.~Lempel, \emph{A universal algorithm for sequential data
compression}, IEEE Trans.\ Inf.\ Theory, vol.~23, 1977.

\bibitem{kaspi1985}
A.~H.~Kaspi, \emph{Two-way source coding with a fidelity criterion},
IEEE Trans.\ Inf.\ Theory, vol.~31, 1985.

\bibitem{orlitsky1990}
A.~Orlitsky, \emph{Worst-case interactive communication I: Two
messages are almost optimal}, IEEE Trans.\ Inf.\ Theory, vol.~36, 1990.

\bibitem{landauer1961}
R.~Landauer, \emph{Irreversibility and heat generation in the
computing process}, IBM J.\ Research and Development, vol.~5, 1961.

\bibitem{tishby1999}
N.~Tishby, F.~C.~Pereira, W.~Bialek, \emph{The information bottleneck
method}, Proc.\ 37th Allerton Conf., 1999.

\bibitem{popovski2022}
P.~Popovski, O.~Simeone, F.~Boccardi, D.~Gündüz, O.~Sahin,
\emph{Semantic-effectiveness filtering and control for post-5G
wireless connectivity}, Journal of the Indian Institute of Science,
2020; and subsequent literature on semantic communication, review
2022.

\end{thebibliography}

\end{document}
